\begin{lemma}
\label{lemma-helper-algebraize}
In Situation \ref{situation-algebraize} let $(M_n)$ be an inverse system of
$A$-modules as in Lemma \ref{lemma-system-of-modules} and let
$(\mathcal{F}_n)$ be the corresponding object of
$\textit{Coh}(U, I\mathcal{O}_U)$. Let $d \geq \text{cd}(A, I)$
and $s \geq 0$ be integers.
With notation as above assume
\begin{enumerate}
\item $A$ is local with maximal ideal $\mathfrak m = \mathfrak a$,
\item $A$ has a dualizing complex, and
\item $(\mathcal{F}_n)$ satisfies the $(s, s + d)$-inequalities
(Definition \ref{definition-s-d-inequalities}).
\end{enumerate}
Let $E$ be an injective hull of the residue field of $A$. Then for $i \leq s$
there exists a finite $A$-module $N$ annihilated by a power
of $I$ and for $n \gg 0$ compatible maps
$$
H^i_\mathfrak m(M_n) \to \Hom_A(N, E)
$$
whose cokernels are finite length $A$-modules and whose kernels $K_n$
form an inverse system such that $\Im(K_{n''} \to K_{n'})$ has finite
length for $n'' \gg n' \gg 0$.
\end{lemma}

\begin{proof}
Let $\omega_A^\bullet$ be a normalized dualizing complex. Then
$\delta^Y_Z = \delta$ is the dimension function associated with
this dualizing complex.
Observe that $\Ext^{-i}_A(M_n, \omega_A^\bullet)$ is a finite $A$-module
annihilated by $I^n$. Fix $0 \leq i \leq s$.
Below we will find $n_1 > n_0 > 0$ such that if we set
$$
N = \Im(\Ext^{-i}_A(M_{n_0}, \omega_A^\bullet) \to
\Ext^{-i}_A(M_{n_1}, \omega_A^\bullet))
$$
then the kernels of the maps
$$
N \to \Ext^{-i}_A(M_n, \omega_A^\bullet),\quad n \geq n_1
$$
are finite length $A$-modules and the cokernels $Q_n$ form a
system such that $\Im(Q_{n'} \to Q_{n''})$ has finite length
for $n'' \gg n' \gg n_1$. This is equivalent to the statement that
the system $\{\Ext^{-i}_A(M_n, \omega_A^\bullet)\}_{n \geq 1}$
is essentially constant in the quotient of the category of finite
$A$-modules modulo the Serre subcategory of finite length $A$-modules.
By the local duality theorem
(Dualizing Complexes, Lemma \ref{dualizing-lemma-special-case-local-duality})
and Matlis duality
(Dualizing Complexes, Proposition \ref{dualizing-proposition-matlis})
we conclude that there are maps
$$
H^i_\mathfrak m(M_n) \to \Hom_A(N, E),\quad n \geq n_1
$$
as in the statement of the lemma.

\medskip\noindent
Pick $f \in \mathfrak m$. Let $B = A_f^\wedge$ be the $I$-adic completion
of the localization $A_f$. Recall that
$\omega_{A_f}^\bullet = \omega_A^\bullet \otimes_A A_f$
and $\omega_B^\bullet = \omega_A^\bullet \otimes_A B$ are dualizing
complexes (Dualizing Complexes, Lemma \ref{dualizing-lemma-dualizing-localize}
and \ref{dualizing-lemma-completion-henselization-dualizing}).
Let $M$ be the finite $B$-module $\lim M_{n, f}$ (compare with
discussion in Cohomology of Schemes, Lemma
\ref{coherent-lemma-inverse-systems-affine}). Then
$$
\Ext^{-i}_A(M_n, \omega_A^\bullet)_f =
\Ext^{-i}_{A_f}(M_{n, f}, \omega_{A_f}^\bullet) =
\Ext^{-i}_B(M/I^n M, \omega_B^\bullet)
$$
Since $\mathfrak m$ can be generated by finitely many $f \in \mathfrak m$
it suffices to show that for each $f$ the system
$$
\{\Ext^{-i}_B(M/I^n M, \omega_B^\bullet)\}_{n \geq 1}
$$
is essentially constant. Some details omitted.

\medskip\noindent
Let $\mathfrak q \subset IB$ be a prime ideal. Then $\mathfrak q$ corresponds
to a point $y \in U \cap Y$. Observe that
$\delta(\mathfrak q) = \dim(\overline{\{y\}})$
is also the value of the dimension function associated to $\omega_B^\bullet$
(we omit the details; use that $\omega_B^\bullet$ is gotten from
$\omega_A^\bullet$ by tensoring up with $B$). Assumption
(3) guarantees via Lemma \ref{lemma-elementary}
that Lemma \ref{lemma-algebraize-local-cohomology-bis}
applies to
$B_\mathfrak q, IB_\mathfrak q, \mathfrak qB_\mathfrak q, M_\mathfrak q$
with $s$ replaced by $s - \delta(y)$. We obtain that
$$
H^{i - \delta(\mathfrak q)}_{\mathfrak qB_\mathfrak q}(M_\mathfrak q) =
\lim H^{i - \delta(\mathfrak q)}_{\mathfrak qB_\mathfrak q}(
(M/I^nM)_\mathfrak q)
$$
and this module is annihilated by a power of $I$.
By Lemma \ref{lemma-terrific} we find that the inverse systems
$H^{i - \delta(\mathfrak q)}_{\mathfrak qB_\mathfrak q}((M/I^nM)_\mathfrak q)$
are essentially constant with value
$H^{i - \delta(\mathfrak q)}_{\mathfrak qB_\mathfrak q}(M_\mathfrak q)$.
Since $(\omega_B^\bullet)_\mathfrak q[-\delta(\mathfrak q)]$ is a normalized
dualizing complex on $B_\mathfrak q$ the local duality theorem
shows that the system
$$
\Ext^{-i}_B(M/I^n M, \omega_B^\bullet)_\mathfrak q
$$
is essentially constant with value
$\Ext^{-i}_B(M, \omega_B^\bullet)_\mathfrak q$.

\medskip\noindent
To finish the proof we globalize as in the proof of
Lemma \ref{lemma-bootstrap}; the argument here is easier
because we know the value of our system already. Namely, consider the maps
$$
\alpha_n :
\Ext^{-i}_B(M/I^n M, \omega_B^\bullet)
\longrightarrow
\Ext^{-i}_B(M, \omega_B^\bullet)
$$
for varying $n$. By the above, for every $\mathfrak q$ we can find an
$n$ such that $\alpha_n$ is surjective after localization at $\mathfrak q$.
Since $B$ is Noetherian and $\Ext^{-i}_B(M, \omega_B^\bullet)$
a finite module, we can find an $n$ such that $\alpha_n$ is surjective.
For any $n$ such that $\alpha_n$ is surjective, given a prime
$\mathfrak q \in V(IB)$ we can find an $n' > n$ such that
$\Ker(\alpha_n)$ maps to zero in $\Ext^{-i}(M/I^{n'}M, \omega_B^\bullet)$
at least after localizing at $\mathfrak q$.
Since $\Ker(\alpha_n)$ is a finite $A$-module and since supports of
sections are quasi-compact, we can find an $n'$ such that
$\Ker(\alpha_n)$ maps to zero in $\Ext^{-i}(M/I^{n'}M, \omega_B^\bullet)$.
In this way we see that $\Ext^{-i}(M/I^n M, \omega_B^\bullet)$
is essentially constant with value $\Ext^{-i}(M, \omega_B^\bullet)$.
This finishes the proof.
\end{proof}
